\section{Tokenomics}
\label{sec:tokenomics}

Every number in this section is taken from the tokenomics record
(\code{TOKENOMICS.md}), which is the single source of truth for the chain's
economic parameters. Where the record and the running code disagree, the code
is what validators execute and the record is corrected; this document is
re-released when a number moves.

\subsection{The \kash{} token}

\begin{table}[H]
\centering
\small
\begin{tabular}{@{}ll@{}}
\toprule
Symbol & \kash{} \\
Base denomination & \esp{}, 18 decimals (1\,\kash{} $= 10^{18}$\,\esp) \\
Gas token & \kash{} --- the same asset; there is no second token \\
Genesis supply & \textbf{1\,000\,000\,000 \kash{}} \\
Maximum supply & uncapped (\code{x/mint} \code{max\_supply = 0}); see \S\ref{sec:netsupply} \\
Uses & \legal{gas, staking and delegation, governance deposits and voting} \\
\bottomrule
\end{tabular}
\caption{The token (TOKENOMICS \S1).}
\end{table}

\subsection{Issuance}
\label{sec:issuance}

New \kash{} is minted every block as a pure function of how much \kash{} is
bonded to validators, modelled on Ethereum's post-merge issuance
(\dref{D4}):

\begin{equation}
\text{annual issuance (\kash)} = F \times \sqrt{\text{bonded \kash}},
\qquad F = 1265
\label{eq:issuance}
\end{equation}
\begin{equation}
\text{issuance rate relative to bonded stake} = \frac{F}{\sqrt{\text{bonded \kash}}}
\label{eq:apr}
\end{equation}

There is \textbf{no} bonded-ratio target and no rate-of-change loop. The Cosmos
SDK's default minter, which steers inflation within a 7--20\,\% band toward a
67\,\% bonded ratio, is switched off: those parameters are zeroed in genesis and
ignored. The curve is implemented as the stock \code{x/mint} module's
\code{MintFn} hook, so it adds no custom module to the chain.

$F$ was chosen against the 1\,B \kash{} genesis supply so that the curve
issues 8\,\% of bonded stake per year at 25\,\% bonded and 4\,\% with
everything bonded:

\begin{table}[H]
\centering
\small
\begin{tabular}{@{}rrrrr@{}}
\toprule
Bonded share & Bonded \kash{} & Issuance $\div$ bonded & Issued / year & Issued as \% of supply \\
\midrule
10\,\% & 100\,M & 12.65\,\% & 12.65\,M & 1.27\,\% \\
25\,\% & 250\,M & \textbf{8.00\,\%} & \textbf{20.0\,M} & \textbf{2.00\,\%} \\
50\,\% & 500\,M & 5.66\,\% & 28.3\,M & 2.83\,\% \\
75\,\% & 750\,M & 4.62\,\% & 34.6\,M & 3.46\,\% \\
100\,\% & 1\,B & 4.00\,\% & 40.0\,M & 4.00\,\% \\
\bottomrule
\end{tabular}
\caption{Issuance against bonded stake, $F = 1265$, 1\,B supply
(TOKENOMICS \S2). At the planned genesis bonding of $\approx$120\,M \kash{}
(\S\ref{sec:allocation}) the curve issues $\approx$11.5\,\% of bonded stake
per year.}
\label{tab:issuance}
\end{table}

Properties of the curve:

\begin{itemize}
\item Doubling bonded stake raises total issuance by $\sqrt{2}$ ($\times 1.41$)
  and cuts issuance per unit of bonded stake by the same factor. Issuance never
  reaches zero and is never negative.
\item Because issuance grows sub-linearly in stake, the inflation rate
  (issued $\div$ supply) is bounded above by $F / \sqrt{\text{supply}}$ ---
  4\,\% of the genesis supply with everything bonded.
\item Issuance is chain-wide and distributed pro-rata by bonded stake, so the
  size of the active validator set does not change how it is divided per unit
  of stake, and the curve adds no concentration incentive of its own.
\item $F$ is a \textbf{protocol constant}, in the sense of Ethereum's
  \code{BASE\_REWARD\_FACTOR}: changing it is a coordinated software upgrade
  reviewed and voted as such, not a governance parameter change.
\item Per-block provision is the annual figure divided by \code{blocks\_per\_year},
  which \emph{is} a governance parameter and must track real block time. It is
  currently set for $\approx$1.5\,s blocks (21\,038\,400) and will be re-derived
  from observed \code{testnet-1} block time before mainnet; if blocks are
  slower than assumed, realised issuance is proportionally lower than
  Table~\ref{tab:issuance}, and vice versa.
\end{itemize}

\subsubsection{Where issued \kash{} goes}

Minted coins enter the fee collector and are swept by the distribution module
every block, together with transaction tips (\S\ref{sec:fees}):

\begin{enumerate}
\item A \textbf{community tax} of \textbf{2\,\%} goes to the on-chain community
  pool, spendable only by governance proposal.
\item The remainder goes to validators pro-rata by bonded stake. Each validator
  takes its commission (chain-wide floor 5\,\%, \dref{D10}) and passes the rest to
  its delegators pro-rata.
\end{enumerate}

Of the issuance attributable to a delegation, the 2\,\% community tax and
the validator's commission are deducted before the remainder is credited.
Issuance credited to stakers is protocol issuance, variable with bonded stake
and block time, and reduced by slashing; it is not a return offered by anyone
(\S\ref{sec:disclaimers}).

\subsection{Fees and burn}
\label{sec:fees}

Konstellation runs EIP-1559 for every transaction, EVM and Cosmos alike
(\dref{D5}).

\begin{table}[H]
\centering
\small
\begin{tabularx}{\textwidth}{@{}llL@{}}
\toprule
Parameter & Value & Note \\
\midrule
Initial base fee & $10^{9}$ \esp{} per gas (1 gwei-equivalent) & \code{cosmos/evm} default; adjusts from block~1 \\
Base-fee change denominator & 8 & at most $\pm$12.5\,\% per block, as Ethereum \\
Elasticity multiplier & 2 & target utilisation is half the block gas limit \\
\code{min\_gas\_price} (base-fee floor) & 0 & no protocol floor; the base fee may decay toward zero when the chain is idle. Governance parameter. The foundation's public RPC nodes apply a node-local minimum of $10^{9}$ \esp{} per gas (1 gwei-equivalent) as a spam guard on what they accept into their mempool; that is a node setting, not a protocol rule, and other nodes may set their own \\
\code{min\_gas\_multiplier} & 0.5 & a block-level floor on the gas-wanted figure the fee market records, preventing a proposer from pushing the base fee down by over-reporting; not a user-facing charge \\
\bottomrule
\end{tabularx}
\caption{Fee-market parameters (TOKENOMICS \S3).}
\end{table}

Every transaction pays $\text{gasUsed} \times (\text{baseFee} + \text{tip})$.
Of that:

\begin{itemize}
\item \textbf{$\text{baseFee} \times \text{gasUsed}$ is burned} --- destroyed at
  the end of every block, before the distribution module can see it. This is
  Ethereum semantics and replaces the Cosmos SDK default of routing the base fee
  to validators.
\item \textbf{The tip} (priority fee) goes to validators and delegators through
  the same pipeline as issuance, including the 2\,\% community tax.
\end{itemize}

For Cosmos-SDK transactions, which prepay for gas \emph{wanted}, the burn is still
computed on gas \emph{used}; the $\text{baseFee} \times (\text{gasWanted} -
\text{gasUsed})$ slack goes to validators like a tip.

\subsection{Net supply}
\label{sec:netsupply}

\begin{equation}
\Delta\,\text{supply per block} = \text{issuance}(\text{bonded}) - \text{baseFee} \times \text{gasUsed}
\end{equation}

Burn scales with usage and congestion; issuance scales with stake. Order of
magnitude, at 1.5\,s blocks ($\approx$21\,M blocks per year):

\begin{table}[H]
\centering
\small
\begin{tabular}{@{}rrr@{}}
\toprule
Average base fee & Average gas per block & Burned per year \\
\midrule
1 gwei & 1\,M & 21\,k \kash{} \\
1 gwei & 10\,M & 210\,k \kash{} \\
10 gwei & 10\,M & 2.1\,M \kash{} \\
50 gwei & 20\,M & 21\,M \kash{} \\
\bottomrule
\end{tabular}
\caption{Burn against usage (TOKENOMICS \S3.1). ``gwei'' here means
$10^{9}$ \esp{}.}
\end{table}

Against $\approx$20\,M \kash{} per year issued at 25\,\% bonded, the chain is
\textbf{net inflationary until sustained demand pushes the base fee into the
tens of gwei at high utilisation} --- the same shape as Ethereum, where net
deflation has only occurred in periods of heavy use. This is expected and is
not a design flaw: burn is a function of demand for block space; issuance is
a function of the security budget.

\subsection{Staking parameters}
\label{sec:staking}

\begin{table}[H]
\centering
\small
\begin{tabularx}{\textwidth}{@{}lL@{}}
\toprule
Consensus & delegated proof of stake (\code{x/staking}) \\
Active set (\code{max\_validators}) & \textbf{30} \\
Unbonding period & \textbf{21 days} \\
Minimum validator commission & \textbf{5\,\%} \\
Double-sign slash & \textbf{5\,\%} of stake; permanent tombstone \\
Downtime slash & \textbf{0.01\,\%} of stake \\
Downtime definition & missing more than 50\,\% of a 100-block window (SDK default) \\
Downtime jail & 10 minutes, then an unjail transaction (SDK default) \\
Launch validator set & 10 foundation-run validators at genesis, further admission permissioned (\dref{D7}, \dref{D16}, \S\ref{sec:validators}); same on \code{testnet-1} and mainnet \\
\bottomrule
\end{tabularx}
\caption{Staking and slashing parameters (\dref{D10}, TOKENOMICS \S4). Slashed
stake is burned.}
\end{table}

The active-set cap of 30 is a deliberate DPoS choice, not an SDK default
carried over; the roadmap for raising it is in \S\ref{sec:validators}.

\subsection{Governance parameters}
\label{sec:governance}

\begin{table}[H]
\centering
\small
\begin{tabularx}{\textwidth}{@{}llL@{}}
\toprule
Parameter & Mainnet & Note \\
\midrule
Minimum proposal deposit & \textbf{1\,000 \kash{}} & to enter the voting period. \code{testnet-1}: 10 \\
Expedited minimum deposit & \textbf{5\,000 \kash{}} & \code{testnet-1}: 50 \\
Maximum deposit period & 2 days & SDK default \\
Voting period & \textbf{3 days} & to be lengthened as the validator set decentralises. \code{testnet-1}: 2 hours \\
Expedited voting period & 1 day & SDK default. \code{testnet-1}: 30 minutes \\
Quorum & \textbf{33.4\,\%} & of bonded stake \\
Pass threshold & \textbf{50\,\%} & of non-abstain votes \\
Expedited pass threshold & 66.7\,\% & SDK default \\
Veto threshold & 33.4\,\% & SDK default \\
Deposit on veto & \textbf{burned} & \\
Deposit on failed quorum, rejection or pass & \textbf{refunded} & \\
\bottomrule
\end{tabularx}
\caption{Governance parameters (\dref{D11}, TOKENOMICS \S5). The mainnet values
are selected only by the exact chain-id \code{konstellation-1}; every other
chain-id, including local development, gets the testnet profile.}
\end{table}

Deposits are refunded on every outcome except veto, so \textbf{the deposit is a
spam bond, not a cost}: it decides who can afford to \emph{propose}, not what
proposing costs. 1\,000 \kash{} is $10^{-6}$ of the genesis supply --- inside
what any serious proposer holds, and well above what a spammer wants locked for
up to five days per proposal. All values in the table are governance parameters
and are adjustable by proposal.

Governance can also override the compliance authority
(\S\ref{sec:compliance}), set and remove IBC rate limits (\S\ref{sec:rails}),
and is the authority of the circuit breaker (\dref{D14}).

\subsection{Genesis allocation}
\label{sec:allocation}

On a 1\,000\,000\,000 \kash{} supply (TOKENOMICS \S7, decided 2026-09-15;
liquid tranches and community split settled 2026-09-20) the allocation is
given in Table~\ref{tab:allocation}, with the community bucket split out in
Table~\ref{tab:community}.

\begin{table}[tbp]
\centering
\small
\begin{tabularx}{\textwidth}{@{}>{\raggedright\arraybackslash}p{2.4cm}rrp{1.9cm}LL@{}}
\toprule
Bucket & Share & \kash{} & \raggedleft Liquid at genesis & Vesting / release & Held by \\
\midrule
Founding team \& early contributors & 22\,\% & 220\,M & \raggedleft 22\,M (10\,\% of each grant) & remaining 90\,\%: 12-month cliff, then linear over 36 months (4 years total) & liquid part paid to each member's address in genesis; locked part in a revocable vesting wallet, one per person (\dref{D12}) \\
Community \& developers & 33\,\% & 330\,M & \raggedleft 50\,M (the community-pool seed) & remaining 280\,M released over 5 years, front-loaded & see Table~\ref{tab:community} \\
Treasury (foundation) & 25\,\% & 250\,M & \raggedleft 50\,M & 20\,\% liquid at genesis; remainder linear over 48 months & foundation multisig ($\geq$3-of-5); vesting wallet for the locked part \\
Validator bootstrap \& staking & 12\,\% & 120\,M & \raggedleft 120\,M, bonded & none --- bonded at genesis via gentx & the foundation, as operator of the ten launch validators (\dref{D7}) \\
Liquidity \& public distribution & 8\,\% & 80\,M & \raggedleft 80\,M & none & foundation, earmarked for \legal{DEX liquidity, market making and any public sale} \\
\midrule
 & 100\,\% & 1\,000\,M & \raggedleft 322\,M (32.2\,\%) & & \\
\bottomrule
\end{tabularx}
\caption{Genesis allocation.}
\label{tab:allocation}
\end{table}

\begin{table}[tbp]
\centering
\small
\begin{tabularx}{\textwidth}{@{}>{\raggedright\arraybackslash}p{2.4cm}rrlLL@{}}
\toprule
Sub-bucket & Share & \kash{} & At genesis & Purpose & Steward and release \\
\midrule
Ecosystem \& developer grants & 18\,\% & 180\,M & locked & builders, integrations, tooling, bug bounties & foundation grants committee; large grants ratified by governance. Non-revocable vesting wallets, beneficiary the foundation multisig, five yearly tranches of 30 / 25 / 20 / 15 / 10\,\% (54 / 45 / 36 / 27 / 18\,M) \\
User \& developer incentives & 10\,\% & 100\,M & locked & usage rewards, airdrops, liquidity mining, hackathon prizes & foundation, programme by programme. Same wallet shape; tranches 30 / 25 / 20 / 15 / 10\,\% (30 / 25 / 20 / 15 / 10\,M) \\
On-chain community pool seed & 5\,\% & 50\,M & \textbf{liquid} & governance-spendable from day one, on top of the 2\,\% tax & the \code{x/distribution} community pool; governance proposals only. Written directly into genesis distribution state, never through a contract \\
\bottomrule
\end{tabularx}
\caption{The community \& developers bucket (33\,\%), split.}
\label{tab:community}
\end{table}

The community pool is a Cosmos \emph{module account}: no key holds it, only a
passed governance proposal moves anything out of it, and the chain refuses EVM
value transfers into it. It is therefore the one part of the community bucket
that cannot sit behind a vesting contract; it is seeded in genesis and grows
by the 2\,\% tax. Grants and incentives, which the foundation spends
programmatically, are the parts that vest.

Why these numbers:

\begin{itemize}
\item \legal{\textbf{Team 22\,\%} is the middle of the Layer-1 range.
  \textbf{10\,\% of each grant is liquid at launch}; the other 90\,\% is
  behind a one-year cliff and vests over the following three years, so 90\,\%
  of the team's allocation is non-transferable in year one.}
\item \textbf{Community 33\,\% is the largest bucket.} Grants and incentives are
  the actual growth lever for a new EVM chain.
\item \textbf{Treasury 25\,\%} covers audits (\dref{D9}), infrastructure, legal,
  \legal{listings} and multi-year runway; 20\,\% is liquid because those costs
  start before launch.
\item \legal{\textbf{Validators 12\,\%, bonded at genesis.} $\approx$120\,M bonded on
  day one puts equation~\ref{eq:apr} at $\approx$11.5\,\% of bonded stake per
  year, falling as more stake is bonded (\S\ref{sec:validators}).}
\item \legal{\textbf{32.2\,\% of supply is transferable at genesis} (bonded
  validator stake 120\,M, treasury tranche 50\,M, liquidity 80\,M,
  community-pool seed 50\,M, team 22\,M); of that, 120\,M is bonded behind a
  21-day unbond and 50\,M moves only by governance.}
\end{itemize}

\begin{table}[H]
\centering
\small
\begin{tabular}{@{}lrrrrrr@{}}
\toprule
End of year & Team & Community & Treasury & Validators & Liquidity & \textbf{Circulating} \\
\midrule
genesis & 22\,M & 50\,M & 50\,M & 120\,M & 80\,M & \textbf{322\,M (32.2\,\%)} \\
1 & 22\,M & 134\,M & 100\,M & 120\,M & 80\,M & \textbf{456\,M (45.6\,\%)} \\
2 & 88\,M & 204\,M & 150\,M & 120\,M & 80\,M & \textbf{642\,M (64.2\,\%)} \\
3 & 154\,M & 260\,M & 200\,M & 120\,M & 80\,M & \textbf{814\,M (81.4\,\%)} \\
4 & 220\,M & 302\,M & 250\,M & 120\,M & 80\,M & \textbf{972\,M (97.2\,\%)} \\
5 & 220\,M & 330\,M & 250\,M & 120\,M & 80\,M & \textbf{1\,000\,M (100\,\%)} \\
\bottomrule
\end{tabular}
\caption{Circulating supply by year, issuance excluded. Community: 50\,M pool
seed at genesis plus 280\,M released 30 / 25 / 20 / 15 / 10\,\% per year.
Team: 22\,M at genesis plus 198\,M vesting linearly from month 12 to month 48
(TOKENOMICS \S7).}
\end{table}

Names and individual amounts within the team bucket are not part of this
document. \todo{Decide whether the whitepaper will disclose named beneficiaries
or aggregate only; nothing is recorded either way.}

\subsection{Vesting}
\label{sec:vesting}

All vesting is implemented in \textbf{Solidity vesting contracts}, not in the
Cosmos SDK's \code{x/auth} vesting accounts (\dref{D12}). The reason is defensive:
a 2026 exploit on another Cosmos EVM chain was attributed to a flaw touching
vesting accounts and balance handling, and SDK vesting accounts add a second,
protocol-level balance model that every balance-touching path must get right.
Contracts are auditable in isolation, and a contest-style Solidity audit is a
good fit for them (\S\ref{sec:audit}).

The contracts are built on the OpenZeppelin v5 vesting-wallet base, hold native
\kash{} only, and are deployed at \code{CREATE2}-deterministic addresses known
before genesis, so the genesis file funds each wallet directly. A vesting year
is 365 days. Three shapes are used:

\begin{description}[style=nextline]
\item[Team grants (revocable).] 10\,\% of each grant is a plain genesis balance
  and never touches a contract. The remaining 90\,\% sits in a per-person
  wallet that releases nothing until the 12-month cliff, then linearly over the
  following 36 months. The foundation multisig may \textbf{revoke} a departing
  member's wallet: what has vested stays with the member, what has not returns
  to the treasury. A member who leaves in month eight walks away with only
  the liquid 10\,\%.
\item[Community grants and incentives (non-revocable).] 280\,M in tranche
  wallets whose beneficiary is the foundation multisig, one tranche per year,
  front-loaded: 30 / 25 / 20 / 15 / 10\,\% (Table~\ref{tab:community}). Each
  year's tranche vests linearly within its year rather than unlocking as a
  step. Nobody can revoke or accelerate them.
\item[Treasury (non-revocable).] 50\,M liquid at genesis; 200\,M in a wallet
  vesting linearly over 48 months to the foundation multisig.
\end{description}

The community-pool seed (50\,M) is not a contract at all: it is a module
account written into the genesis distribution state and spendable only by
governance. Calling \code{release()} on any wallet moves whatever has vested so
far to its beneficiary; nobody else can move it.

The vesting contracts are built and pending review at the time of this draft
(in the \code{contracts} repository, \code{src/vesting/}); their audit is part
of the pre-mainnet scope.
\todo{Confirm the final vesting contract parameters once reviewed and merged,
and the signer set of the foundation multisig that holds the team-grant
revocation role (the multisig is decided; its signers are not).}

\subsection{Flow of value, end to end}

\begin{figure}[H]
\footnotesize
\begin{verbatim}
                    +------------- issuance: F x sqrt(bonded) --------------+
                    |                                                       v
users --gas x (base+tip)--> fee collector --- EndBlock: burn base x gasUsed --> (burned)
                                  |
                                  +---- next BeginBlock: x/distribution
                                            |-- 2 % community pool  (gov-spendable)
                                            +-- 98 % validators, pro-rata by stake
                                                  |-- commission >= 5 %  -> validator
                                                  +-- remainder          -> delegators
\end{verbatim}
\caption{Flow of value (TOKENOMICS \S6).}
\end{figure}
